% =============================================================================
% abstrak.tex — Abstrak (Bahasa Indonesia)
% Spasi 1 (single spacing), TNR 12pt, 250-500 kata
% =============================================================================

\chapter*{ABSTRAK}
\addcontentsline{toc}{chapter}{ABSTRAK}
\thispagestyle{plain}

\begin{abstrakenv}

% Baris identitas
\noindent
\textbf{\NamaMahasiswa} (\NIM). \textbf{\NamaProdi}.
\textit{\JudulTA}. Pembimbing: \NamaPembimbing\ dan \NamaPembimbingDua.

\vspace{0.5cm}

% Isi abstrak (250-500 kata)
Tulis abstrak dalam Bahasa Indonesia di sini. Abstrak harus menguraikan secara singkat
latar belakang, tujuan, metode penelitian, hasil, dan kesimpulan dari penelitian.
Panjang abstrak antara 250 sampai 500 kata.

Lorem ipsum dolor sit amet, consectetur adipiscing elit. Sed do eiusmod tempor incididunt
ut labore et dolore magna aliqua. Ut enim ad minim veniam, quis nostrud exercitation
ullamco laboris nisi ut aliquip ex ea commodo consequat. Duis aute irure dolor in
reprehenderit in voluptate velit esse cillum dolore eu fugiat nulla pariatur.
Excepteur sint occaecat cupidatat non proident, sunt in culpa qui officia deserunt
mollit anim id est laborum.

\vspace{0.5cm}

\noindent
\textbf{Kata Kunci:} kata kunci pertama, kata kunci kedua, kata kunci ketiga, kata kunci keempat, kata kunci kelima.

\end{abstrakenv}
